\documentclass{beamer}
\usepackage[T1]{fontenc}
\usepackage{fontspec}
\usepackage{amsmath, amssymb, amsthm}
\usepackage{xcolor, graphicx, booktabs}
\usepackage{tikz}
\usetikzlibrary{positioning, arrows.meta}

% -------------------- Theme settings --------------------
\usetheme{Berlin}
\usecolortheme{whale}
\setmainfont{Times New Roman}

% -------------------- Theorem environments --------------------
\theoremstyle{definition}
\newtheorem{defn}{Definition}
\newtheorem{ex}[defn]{Example}
\newtheorem{remark}[defn]{Remark}
\theoremstyle{plain}
\newtheorem{thm}[defn]{Theorem}
\newtheorem{lem}[defn]{Lemma}
\newtheorem{cor}[defn]{Corollary}
\newtheorem{prop}[defn]{Proposition}

% -------------------- Title page --------------------
\title{Deformation Theory of Calabi-Yau Varieties}
\subtitle{Undergraduate Thesis Proposal}
\author{Weining Huang}
\institute{Beihang University \\ School of Mathematical Sciences}
\date{\today}

\begin{document}

\begin{frame}
    \thispagestyle{empty}
    \maketitle
\end{frame}

\begin{frame}{Outline}
    \tableofcontents
\end{frame}

% ------------------------------------------------------------------------------------------
%     Part 1: Background and Preliminaries
% ------------------------------------------------------------------------------------------
\section{Preliminaries}

\begin{frame}{Convention}
    \begin{itemize}
        \item Usually `Variety' means a seperated integral scheme of finite type over a fixed field $ k $. Here we can think $ k=\mathbb{C}  $. 
        \item But for convenience we assume all varieties appear in the slides are projective.
        \item Thus we can regard varieties over $ \mathbb{C} $ as complex compact K\"{a}hler manifolds and use tools in both algebraic geometry and complex geometry to investigate by GAGA.
    \end{itemize}
\end{frame}

\begin{frame}{Deformation Theory with Complex-geometric Approach}
    For a smooth complex manifold $ X $, we have:
    \pause
    \begin{itemize}
        \item First order deformations of complex structure (of $ X $) are described by \textbf{Kodaira-Spencer class} in $ H^1(T_X) $.
        \pause
        \item Deformations are governed by the \textbf{Maurer-Cartan equation}:
        \[
        \bar{\partial} \varphi + \frac{1}{2}[\varphi, \varphi] = 0, \quad \varphi \in \mathcal{A}^{0,1}(T_X)
        \]
        % \pause
        % \item \textbf{Kuranishi theorem}: For compact $ X $ there exists a versal deformation space $(S,0)$ with tangent space $H^1(T_X)$.
    \end{itemize}
\end{frame}

\begin{frame}{Deformation Theory with Algebraic Approach}
    Let $ X $ be a scheme over a fixed field $ k(=\mathbb{C} ) $.
    \begin{itemize}
        \item Define the deformation functor of $ X $:
        \[
        \text{Def}_X:= \mathsf{C} \to \mathsf{Set},\ A \mapsto \{\text{deformations of}\ X\ \text{over}\ A\}.
        \]Here $ \mathsf{C} $ is he category of local artinian $ k $-algebras with residue field $ k $.
        \pause
        \item \textbf{Schlessinger's Criterion} helps to investigate properties of deformation functor.
    \end{itemize}
\end{frame}

% \subsection{Local Deformation Theory: ODP and Cusp Cases}

%         1111111  odp cusp统一到特定维数

% \begin{frame}{ODP:Definition and 3-fold Geometry}
%     \begin{defn}[Original Double Point]
%         Isolated hypersurface singularity locally $\cong \{ \sum_{i=1}^{4} z_i^2 = 0 \} \subset \mathbb{C}^{4}$.
%     \end{defn}
%     \pause
%     \begin{itemize}
%         \item \textbf{Resolution (dim 3)}: Blow-up replaces point by $Q \cong \mathbb{P}^1\times\mathbb{P}^1$.
%         \pause
%         \item On $Q$, curve classes $A := [\mathbb{P}^1\times\mathrm{pt}]$, $B := [\mathrm{pt}\times\mathbb{P}^1]$.
%         \pause
%         \item Their difference $A-B$ is central in Friedman's condition.
%     \end{itemize}
% \end{frame}

% \begin{frame}{Cusp Singularities: Definition and Smoothing}
%     \begin{defn}[Cusp (surface singularity)]
%         Minimal resolution consists of a cycle of smooth rational curves meeting transversely.
%     \end{defn}
%     \pause
%     \begin{itemize}
%         \item \textbf{Smoothing criterion}: A cusp is smoothable iff minimal resolution of its dual cusp embeds as an anticanonical divisor in a smooth rational surface.
%         \pause
%         % 444444444444 surface的cusp替换掉一部分主要内容
%         \item \textbf{Necessity}: proved by Looijenga in 1981 and Looijenga conjectured the sufficiency holds.
%         \pause
%         \item \textbf{Sufficiency}: proved by Gross-Hacking-Keel and Philip Engel independently.
%     \end{itemize}
% \end{frame}

% rigidity of some varieties
% \begin{frame}{Rigidity of Homogeneous Varieties}
%     \begin{itemize}
%         \item \textbf{Homogeneous varieties}: $\P^n$, Grassmannians, etc.
%         \pause
%         \item \textbf{Borel-Weil-Bott theorem}: $H^i(\P^n, T_{\P^n}) = 0$ for $i>0$
%         \pause
%         \item \textbf{Consequence}: Projective spaces and Grassmannians are \textbf{rigid}
%         \[
%         \operatorname{Def}(\P^n) = \{\text{point}\}
%         \]
%         \pause
%         \item Intuition: these varieties have too much symmetry to admit non-trivial deformations.
%     \end{itemize}
% \end{frame}

% ------------------------------------------------------------------------------------------
%     Part 2: Deformation of Calabi-Yau Varieties
% ------------------------------------------------------------------------------------------

\section{Smooth Calabi-Yau Manifolds}

\begin{frame}{Deformations of Smooth CY Manifolds}
    \begin{defn}
        A compact Kähler manifold $X$ is called a \textbf{Calabi-Yau manifold} if $K_X \cong \mathcal{O}_X$ (trivial canonical bundle).
    \end{defn}
    \pause
    \begin{thm}[Bogomolov-Tian-Todorov, 1980s]
        Deformations of a smooth Calabi-Yau manifold are unobstructed; i.e., the Kuranishi space is smooth.
    \end{thm}
    \pause
    \begin{itemize}
        \item Complex-geometric proof: Calabi-Yau theorem + Tian-Todorov lemma.
    \end{itemize}
\end{frame}

\begin{frame}{Algebraic proof for BTT Theorem}
    \begin{thm}[Z. Ran, Y. Kawamata, 1992]
        Let $ X $ be a scheme over a fixed field $ k(=\mathbb{C} ) $. If the deformation functor $ \text{Def}_X $ satisfies the $T^1$ lifting property, then deformations of $ X $ are unobstructed.
    \end{thm}
    \pause
    \begin{itemize}
        \item It can be proved that for $ X $ being a nonsingular Calabi-Yau variety, $ \text{Def}_X $ satisfies the $ T^1 $ lifting property.
        \pause
        \item Ziv Ran developed the purely algebraic approach, and Yujiro Kawamata modified the $T^1$ lifting property to more natural version.
    \end{itemize}
\end{frame}

% 222222222222222表述修改，顺序修改，friedman往前放

\section{Threefold Calabi-Yau Varities with ODP}

\begin{frame}{Deformations of Threefold CY varieties with ODP}
    Let's shift our scope to singular and threefold case. Usually our main aim is to deform a singular variety to a nonsingular one.
    \begin{defn}[Original Double Point]
        An \textbf{ODP} is an isolated singularity in a threefold locally analytically isomorphic to the origin of 
        \[
        \{\, (z_1,\dots,z_{4})\in\mathbb{C}^{4} \mid z_1^2+z_2^2+\cdots+z_{4}^2 = 0 \,\}.
        \]
    \end{defn}
\end{frame}

\begin{frame}{Deformations of Threefold CY varieties with ODP}
    Let $ X_0 $ be a threefold Calabi-Yau variety in the sense that dualizing sheaf $ \omega_{X_0}\cong \mathcal{O}_{X_0} $, with only ODP as singularities. Let $ X\rightarrow X_0 $ be a small resolution, with exceptional curves $ C_1,\dots,C_r $, and Let $ [C_i] $ be the class of $ C_i $ in $ H^4(X,\mathbb{C} ) $.
    \pause
    \begin{thm}[Friedman, 1986]
        $ X_0 $ is smoothable if
        \begin{itemize}
            \item $ [C_1],[C_r] $ generate $ H^2(X,\varOmega_{X}^2) $
            \item There exists a relation $ a_1[C_1]+\cdots+a_r[C_r]=0 \in H^4(X,\mathbb{C} ) $ with $ a_i\neq 0 $ for every $ i $.
        \end{itemize}
    \end{thm}
    $ X_0 $ is smoothable means $ X_0 $ can be deformed to a nonsingular variety $ X_t $.
\end{frame}

\begin{frame}{Deformations of Threefold CY varieties with ODP}
    \begin{itemize}
        \item Subsequently, independent work of Y. Kawamata, Z. Ran and G. Tian showed that the first assumption is not necessary provided that there exists a resolution $ X'\rightarrow X_0 $ where $ X' $ is K\"{a}hler.
        \item The key point in removing this assumption is a generalization of the unobstructedness theorem for smooth Calabi-Yau manifolds, in any dimension:
    \end{itemize}
    \pause
    \begin{thm}[Kawamata, Ran, Tian]
        Let $ X_0 $ be a Calabi-Yau variety in the sense that dualizing sheaf $ \omega_{X_0}\cong \mathcal{O}_{X_0} $, with only ODP as singularities and there exists a resolution $ X'\rightarrow X_0 $ where $ X' $ is K\"{a}hler. Then deformations of $ X_0 $ are unobstructed.
    \end{thm}
\end{frame}

% \begin{frame}{Topological Interpretation of Friedman's Condition}
%     \begin{itemize}
%         \item The condition is equivalent to the existence of a 3-chain $\overline{\Delta}$ with boundary $\sum \delta_i (A_i-B_i)$.
%         \pause
%         \item Projecting to $X$ yields a 3-cycle $\Delta \in H_3(X,\mathbb{R})$.
%         \pause
%         \item In the smoothing family, this $\Delta$ corresponds to a \textbf{vanishing cycle}.
%         \pause
%         \item Smoothing exists $\Longleftrightarrow$ there exists a non-zero 3-cycle intersecting all singularities.
%     \end{itemize}
%     \begin{center}
%         \begin{tikzpicture}
%             \draw[thick] (0,0) circle (1);
%             \node at (0,0) {$p_i$};
%             \draw[thick, -{Latex[length=3mm]}] (1.5,0) -- (3,0);
%             \node at (2.25,0.3) {blow up};
%             \draw[thick] (4,-1) .. controls (5,0.5) and (6,-0.5) .. (7,0);
%             \node at (5.5,0.8) {$Q_i\cong\P^1\times\P^1$};
%         \end{tikzpicture}
%     \end{center}
% \end{frame}

% \begin{frame}{Higher-Dimensional Generalisation: Rollenske-Thomas}
%     \begin{thm}[Rollenske-Thomas, 2009]
%         Let $X$ be an odd-dimensional nodal Calabi-Yau variety with $n=2m+1\ge 5$. A necessary condition for the existence of a first-order smoothing is a \textbf{non-linear topological condition}:
%         \[
%         \text{a certain vanishing condition involving the Yukawa coupling}
%         \]
%     \end{thm}
%     \pause
%     \begin{itemize}
%         \item Yukawa coupling: $H^{m,m}(Q) \times H^{m,m}(Q) \times H^{m,m}(Q) \to \mathbb{C}$.
%         \item For threefolds ($m=1$) this reduces to Friedman's linear condition.
%     \end{itemize}
% \end{frame}

% \begin{frame}{Recent Developments: Friedman-Laza's Unified Theory}
%     \begin{block}{Friedman-Laza (2024)}
%         Introduced the notion of \textbf{$k$-liminal singularities}, unifying all the above results.
%         \begin{itemize}
%             \item ODP in threefolds is $1$-liminal.
%             \item For $k$-liminal singularities, the smoothing condition becomes a \textbf{$k$-linear or non-linear topological condition}.
%             \item Gives a complete description for weighted homogeneous singularities.
%         \end{itemize}
%     \end{block}
% \end{frame}

% ------------------------------------------------------------------------------------------
%     Part 4: Research Plan and Schedule
% ------------------------------------------------------------------------------------------
\section{Expected Outcomes and References}

\begin{frame}{Expected Outcomes}
    \begin{itemize}
        \item Understand the basic framework of deformation theory.
        \pause
        \item Master the key results in deformation theory of both smooth and singular Calabi-Yau varieties.
    \end{itemize}
\end{frame}

% ------------------------------------------------------------------------------------------
%     References
% ------------------------------------------------------------------------------------------
\begin{frame}[allowframebreaks]{References}
    \begin{thebibliography}{99}
        \bibitem{Fr86} R. Friedman, \textit{Simultaneous resolution of threefold double points}, Math. Ann. 274 (1986), 671--689.
        \bibitem{TT89} A. Todorov, \textit{The Weil-Petersson geometry of the moduli space of $\mathrm{SU}(n\ge 3)$ (Calabi-Yau) manifolds I}, Comm. Math. Phys. 126 (1989), 325--346.
        \bibitem{Tian87} G. Tian, \textit{Smoothness of the universal deformation space of compact Calabi-Yau manifolds and its Petersson-Weil metric}, in Mathematical Aspects of String Theory, 629--646, World Scientific, 1987.
        \bibitem{Ran92} Z. Ran, \textit{Deformations of manifolds with torsion or negative canonical bundle}, J. Algebraic Geom. 1 (1992), 279--291.
        \bibitem{RT09} S. Rollenske, R. Thomas, \textit{Smoothing nodal Calabi-Yau $n$-folds}, J. Topol. 2 (2009), 405--421.
        \bibitem{FL24} R. Friedman, R. Laza, \textit{Deformations of Calabi-Yau varieties with k-liminal singularities}, Forum Math. Sigma (2024).
        \bibitem{Har77} R. Hartshorne, \textit{Algebraic Geometry}, Graduate Texts in Mathematics 52, Springer-Verlag, New York, 1977.
        \bibitem{Huy05} D. Huybrechts, \textit{Complex Geometry: An Introduction}, Universitext, Springer, Berlin-New York, 2005.
        \bibitem{Har10} R. Hartshorne, \textit{Deformation Theory}, Graduate Texts in Mathematics 257, Springer, New York, 2010.
        \bibitem{GJ03} M. Gross, D. Joyce, and D. Huybrechts, \textit{Calabi-Yau Manifolds and Related Geometries}, Universitext, Springer, Berlin, 2003.
    \end{thebibliography}
\end{frame}

\begin{frame}
    \centering
    \Huge{Thank you!}
    
    \vspace{1cm}
    \large{Questions and comments are welcome.}
\end{frame}

\end{document}